\section{合成数据与 Action Calling：构建可扩展训练飞轮}

\subsection{从感知任务到动作任务的迁移}
讲者强调，仅靠视觉理解（看懂界面）不足以形成强 Agent，必须把能力迁移到 ``调用动作''。因此课程引入合成数据管线，把图像问答、OCR、定位、流程推理转化为动作可执行样本。

\begin{importantbox}{Action Calling 的本质}
Action Calling 不是额外功能，而是把模型输出空间从 ``描述'' 切换到 ``执行''。这要求训练数据同时包含意图、推理路径、动作格式与执行反馈。
\end{importantbox}

\begin{figure}[H]
\centering
\includegraphics[width=0.88\textwidth]{frame-06-synthetic-data.jpg}
\caption{视频约 01:02:32：合成数据管线把视觉问题、推理链与动作调用统一到同一训练流程}
\end{figure}

\subsection{典型合成流程}
\begin{enumerate}
    \item 采样任务场景与界面状态；
    \item 构造问题与目标（含定位、读取、操作、验证）；
    \item 使用模型生成候选推理与动作序列；
    \item 借助执行器和评测器筛选高质量样本；
    \item 进入分阶段训练并持续回流失败样本。
\end{enumerate}

\begin{knowledgebox}{为什么分阶段训练有效}
课程中多次提到 ``先学看，再学做，再学长期策略''。如果一开始就端到端学习长链路动作，模型更易陷入高方差和局部最优。分阶段可降低优化难度并提升样本利用率。
\end{knowledgebox}

\subsection{合成数据的边界}
合成样本虽能扩规模，但分布偏差不可避免。课程建议保持真实轨迹与合成轨迹的比例平衡，并持续做线上任务回归，避免模型只在合成分布内表现良好。

\begin{warningbox}{合成数据三大风险}
\begin{itemize}
    \item 语义漂移：目标描述与真实任务意图不一致；
    \item 动作幻觉：生成看似合理但不可执行的操作；
    \item 偏差累积：模型生成的数据反过来强化自身错误模式。
\end{itemize}
\end{warningbox}

\subsection{本章小结}
合成数据的正确用法是 ``补充与加速''，不是替代真实交互数据。课程给出的关键经验是：把执行验证作为合成流程的硬约束。

